\chapter{先知書总结}
\section{地名}
\subsection{提哥亞 (Tekoa)}
\begin{wrapfigure}{r}{0.6\textwidth}
\centering
\includegraphics[width=0.6\textwidth]{ZongJie/figs/TiGeYa1}
\end{wrapfigure}

猶大的一座城，大約位於伯利恆以南10公里，是先知阿摩司的家鄉（摩1:1）。在馬加比和羅馬時代，提哥亞城仍為人所知，今天它的名稱是「塔奎亞」廢墟（Khirbet Taqua；另譯作「提科」廢墟），保存了古名。今日這是一條方圓5英畝的荒廢村落，鮮有學者挖掘研究。在附近曾發現鐵器時代的墳墓。「提哥亞」的意思是「堡壘」、「搭起帳棚」、「號筒」或「戳刺」。

當統帥約押知道大衛王「心裡想念押沙龍」時，便派人前往提哥亞，請一個聰明的婦人來使大衛和押沙龍和好（撒下14:1-20起）。

王國分裂時期，猶大王羅波安一世曾堅固這城（代下11:6）。後來，當猶大王約沙法受到亞捫人、摩押人和西珥山人攻擊的時候，他在「提哥亞的曠野」與民商議對策（代下20:20）。猶大人被擄回歸之後，他們和提哥亞人重建居所（尼3:5, 3:27）。面對敵人圍攻時，先知耶利米亦勸民在提哥亞吹角（耶6:1）。

堡壘 Fort，搭起帳棚 Pitching of tent，號角 Trumpet，剌，插入 To thrust
●Tequ，位於伯利恆之南約 9 公里，希伯崙東北約 15 公里，海拔高有 822 公尺，城的南北各有一條巨大的深溝，故視野開闊，可望見橄欖山和尼波山。提哥亞曠野是在城的四周 ，古時生產羊群和橄欖，很是富庶，山區中多山洞，有利於畜牧。現今設有為猿猴和赤腳人的保護區。

撒下14:2 約押使一個提哥亞的婦人去說服大衛王，讓押沙龍能回到耶路撒冷。

撒下23:26，代上27:9 大衛的勇士以拉是提哥亞人。

代下11:5 猶大王羅波安為保障，所修築的十四個堅固城之一。

代下20:20 猶大王約沙法與亞捫人，摩押人，米烏利人和西珥山人在提哥亞的曠野會戰，結果約沙法大勝。

尼3:5 提哥亞人參加建立耶路撒冷城牆魚門的部份工作。

耶6:1 先知預言猶大之災禍說，便雅憫人哪，你們要逃出耶路撒冷，在提哥亞吹角，因為有災禍與大毀滅。

摩1:1 大地震前二年，提哥亞牧人中的阿摩斯得默示，對以色列等列邦作預言。

Tekoite）〔屬於提哥亞〕

提哥亞居民。（代下11:6；耶6:1）這個名稱用來指益吉，大衛手下的勇士伊拉的父親（撒下23:26；代上11:28；27:9）；也用來指一個聰明的婦人，她在約押的吩咐下，喬裝成寡婦進見大衛，施計令王收回成命，讓被逐的押沙龍回來（撒下14:2,4,9）。以色列人從巴比倫獲釋回鄉後，提哥亞人有份修復耶路撒冷城牆，雖然他們的「權貴」（和合翻作「貴胄」）不肯效勞。（尼3:5,27）

\subsection{到那日}
\subsubsection{以賽亞書}
以賽亞書 2:2-4/10-22
\textcolor{red}{末後的日子}，耶和華殿的山必堅立超乎諸山，高舉過於萬嶺，萬民都要流歸這山。 3 必有許多國的民前往，說：「來吧！我們登耶和華的山，奔雅各神的殿。主必將他的道教訓我們，我們也要行他的路，因為訓誨必出於錫安，耶和華的言語必出於耶路撒冷。」 4 他必在列國中施行審判，為許多國民斷定是非。他們要將刀打成犁頭，把槍打成鐮刀。這國不舉刀攻擊那國，他們也不再學習戰事。
你當進入巖穴，藏在土中，躲避耶和華的驚嚇和他威嚴的榮光。\textcolor{red}{到那日}，眼目高傲的必降為卑，性情狂傲的都必屈膝，唯獨耶和華被尊崇。必有萬軍耶和華降罰的一個日子，要臨到驕傲狂妄的，一切自高的都必降為卑； 13 又臨到黎巴嫩高大的香柏樹和巴珊的橡樹； 14 又臨到一切高山的峻嶺； 15 又臨到高臺和堅固城牆； 16 又臨到他施的船隻，並一切可愛的美物。 17 驕傲的必屈膝，狂妄的必降卑，在那日，唯獨耶和華被尊崇。偶像必全然廢棄。 19 耶和華興起使地大震動的時候，人就進入石洞，進入土穴，躲避耶和華的驚嚇和他威嚴的榮光。\textcolor{red}{到那日}，人必將為拜而造的金偶像、銀偶像拋給田鼠和蝙蝠。 21 到耶和華興起使地大震動的時候，人好進入磐石洞中和巖石穴裡，躲避耶和華的驚嚇和他威嚴的榮光。 22 你們休要倚靠世人，他鼻孔裡不過有氣息，他在一切事上可算什麼呢？

以賽亞書 5:13-17/24-30
所以我的百姓因無知就被擄去，他們的尊貴人甚是飢餓，群眾極其乾渴。 14 故此，陰間擴張其欲，開了無限量的口，他們的榮耀、群眾、繁華並快樂的人，都落在其中。 15 卑賤人被壓服，尊貴人降為卑，眼目高傲的人也降為卑。 16 唯有萬軍之耶和華因公平而崇高，聖者神因公義顯為聖。\textcolor{red}{那時}，羊羔必來吃草，如同在自己的草場，豐肥人的荒場被遊行的人吃盡。
火苗怎樣吞滅碎秸，乾草怎樣落在火焰之中，照樣，他們的根必像朽物，他們的花必像灰塵飛騰。因為他們厭棄萬軍之耶和華的訓誨，藐視以色列聖者的言語。 25 所以耶和華的怒氣向他的百姓發作，他的手伸出攻擊他們，山嶺就震動，他們的屍首在街市上好像糞土。雖然如此，他的怒氣還未轉消，他的手仍伸不縮。

26 他必豎立大旗招遠方的國民，發嘶聲叫他們從地極而來。看哪，他們必急速奔來！ 27 其中沒有疲倦的、絆跌的，沒有打盹的、睡覺的，腰帶並不放鬆，鞋帶也不折斷。 28 他們的箭快利，弓也上了弦；馬蹄算如堅石，車輪好像旋風。 29 他們要吼叫像母獅子，咆哮像少壯獅子；他們要咆哮抓食，坦然叼去，無人救回。\textcolor{red}{那日}，他們要向以色列人吼叫，像海浪砰訇。人若望地，只見黑暗艱難，光明在雲中變為昏暗。

以賽亞書 8:20-22
人當以訓誨和法度為標準。」他們所說的若不與此相符，必不得見晨光。 21 他們必經過這地，受艱難，受飢餓，飢餓的時候心中焦躁，咒罵自己的君王和自己的神。 22 仰觀上天，俯察下地，不料，盡是艱難黑暗和幽暗的痛苦，他們必被趕入烏黑的黑暗中去。

以賽亞書 10:20-32
\textcolor{red}{到那日}，以色列所剩下的和雅各家所逃脫的，不再倚靠那擊打他們的，卻要誠實倚靠耶和華以色列的聖者。 21 所剩下的，就是雅各家所剩下的，必歸回全能的神。 22 以色列啊，你的百姓雖多如海沙，唯有剩下的歸回。原來滅絕的事已定，必有公義施行，如水漲溢。 23 因為主萬軍之耶和華在全地之中，必成就所定規的結局。所以主萬軍之耶和華如此說：「住錫安我的百姓啊，亞述王雖然用棍擊打你，又照埃及的樣子舉杖攻擊你，你卻不要怕他。 25 因為還有一點點時候，向你們發的憤恨就要完畢，我的怒氣要向他發作，使他滅亡。」 26 萬軍之耶和華要興起鞭來攻擊他，好像在俄立磐石那裡殺戮米甸人一樣；耶和華的杖要向海伸出，把杖舉起，像在埃及一樣。\textcolor{red}{到那日}，亞述王的重擔必離開你的肩頭，他的軛必離開你的頸項，那軛也必因肥壯的緣故撐斷(因膏油的緣故毀壞)。28 亞述王來到亞葉，經過米磯崙，在密抹安放輜重。 29 他們過了隘口，在迦巴住宿。拉瑪人戰兢，掃羅的基比亞人逃跑。 30 迦琳的居民(女子)哪，要高聲呼喊！萊煞人哪，須聽！哀哉，困苦的亞拿突啊！ 31 瑪得米那人躲避，基柄的居民逃遁。\textcolor{red}{當那日}，亞述王要在挪伯歇兵，向錫安女子的山，就是耶路撒冷的山，掄手攻他。

以賽亞書 11:10-16
\textcolor{red}{到那日}，耶西的根立做萬民的大旗，外邦人必尋求他，他安息之所大有榮耀。\textcolor{red}{當那日}，主必二次伸手救回自己百姓中所餘剩的，就是在亞述、埃及、巴忒羅、古實、以攔、示拿、哈馬並眾海島所剩下的。 12 他必向列國豎立大旗，招回以色列被趕散的人，又從地的四方聚集分散的猶大人。 13 以法蓮的嫉妒就必消散，擾害猶大的必被剪除；以法蓮必不嫉妒猶大，猶大也不擾害以法蓮。 14 他們要向西飛，撲在非利士人的肩頭上(西界)，一同擄掠東方人，伸手按住以東和摩押，亞捫人也必順服他們。 15 耶和華必使埃及海汊枯乾，掄手用暴熱的風使大河分為七條，令人過去不致濕腳。 16 為主餘剩的百姓，就是從亞述剩下回來的，必有一條大道，\textcolor{red}{如當日}以色列從埃及地上來一樣。

以賽亞書 19:21-25
耶和華必被埃及人所認識，\textcolor{red}{在那日}，埃及人必認識耶和華，也要獻祭物和供物敬拜他，並向耶和華許願還願。耶和華必擊打埃及，又擊打又醫治，埃及人就歸向耶和華。他必應允他們的禱告，醫治他們。\textcolor{red}{當那日}，必有從埃及通亞述去的大道，亞述人要進入埃及，埃及人也進入亞述，埃及人要與亞述人一同敬拜耶和華。\textcolor{red}{當那日}，以色列必與埃及、亞述三國一律，使地上的人得福。因為萬軍之耶和華賜福給他們，說：「埃及我的百姓，亞述我手的工作，以色列我的產業，都有福了！」


以賽亞書 26:1-7
\textcolor{red}{當那日}，在猶大地人必唱這歌說：「我們有堅固的城，耶和華要將救恩定為城牆，為外郭。 2 敞開城門，使守信的義民得以進入！ 3 堅心倚賴你的，你必保守他十分平安，因為他倚靠你。 4 你們當倚靠耶和華，直到永遠，因為耶和華是永久的磐石。 5 他使住高處的與高城一併敗落，將城拆毀，拆平直到塵埃。 6 要被腳踐踏，就是被困苦人的腳和窮乏人的腳踐踏。」 7 義人的道是正直的，你為正直的主，必修平義人的路。

以賽亞書 27:1-3,12,13
\textcolor{red}{到那日}，耶和華必用他剛硬有力的大刀刑罰鱷魚，就是那快行的蛇，刑罰鱷魚，就是那曲行的蛇，並殺海中的大魚\textcolor{red}{ 當那日}，有出酒的葡萄園，你們要指這園唱歌說： 3 「我耶和華是看守葡萄園的，我必時刻澆灌，晝夜看守，免得有人損害。 4 我心中不存憤怒，唯願荊棘蒺藜與我交戰，我就勇往直前，把她一同焚燒。 5 不然，讓它持住我的能力，使她與我和好；願她與我和好！」 6 將來雅各要扎根，以色列要發芽開花，他們的果實必充滿世界。主擊打他們，豈像擊打那些擊打他們的人嗎？他們被殺戮，豈像被他們所殺戮的嗎？ 8 你打發他們去，是相機宜與他們相爭，颳東風的日子，就用暴風將他們逐去。 9 所以雅各的罪孽得赦免，他的罪過得除掉的果效，全在乎此，就是他叫祭壇的石頭變為打碎的灰石，以致木偶和日像不再立起。 10 因為堅固城變為淒涼，成了撇下離棄的居所，像曠野一樣。牛犢必在那裡吃草，在那裡躺臥，並吃盡其中的樹枝。 11 枝條枯乾，必被折斷，婦女要來點火燒著。因為這百姓蒙昧無知，所以創造他們的必不憐恤他們，造成他們的也不施恩於他們。以色列人哪，\textcolor{red}{到那日}，耶和華必從大河直到埃及小河，將你們一一地收集，如同人打樹拾果一樣。\textcolor{red}{當那日}，必大發角聲，在亞述地將要滅亡的，並在埃及地被趕散的都要來，他們就在耶路撒冷聖山上敬拜耶和華。
馬太福音 24:31
他要差遣使者，用號筒的大聲，將他的選民從四方，從天這邊到天那邊，都招聚了來。
利未記 25:9
當年七月初十日，你要大發角聲，這日就是贖罪日，要在遍地發出角聲。



以賽亞書 56:8
主耶和華，就是招聚以色列被趕散的，說：「在這被招聚的人以外，我還要招聚別人歸併他們。」


\subsubsection{耶利米書}
耶利米書 3:17
\textcolor{red}{那時}，人必稱耶路撒冷為耶和華的寶座，萬國必到耶路撒冷，在耶和華立名的地方聚集。他們必不再隨從自己頑梗的惡心行事。

耶利米書 23:7,8
耶和華說：「\textcolor{red}{日子}將到，人必不再指著那領以色列人從埃及地上來永生的耶和華起誓，
耶利米書 23:20…
耶和華的怒氣必不轉消，直到他心中所擬定的成就了，\textcolor{red}{末後的日子}你們要全然明白。


耶利米書 30:8-11
萬軍之耶和華說：「\textcolor{red}{到那日}，我必從你頸項上折斷仇敵的軛，扭開他的繩索，外邦人不得再使你做他們的奴僕。

耶利米書 30:24
耶和華的烈怒必不轉消，直到他心中所擬定的成就了，\textcolor{red}{末後的日子}你們要明白。

耶利米書 48:47
耶和華說：「\textcolor{red}{到末後}，我還要使被擄的摩押人歸回。」摩押受審判的話到此為止。

耶利米書 49:39
「\textcolor{red}{到末後}，我還要使被擄的以攔人歸回。這是耶和華說的。」


\subsubsection{以西結書}
以西結書 20:34
我必用大能的手和伸出來的膀臂，並傾出來的憤怒，將你們從萬民中領出來，從分散的列國內聚集你們。

以西結書 20:41
我從萬民中領你們出來，從分散的列國內聚集你們，那時我必悅納你們好像馨香之祭，要在外邦人眼前，在你們身上顯為聖。

以西結書 38:16
歌革啊，你必上來攻擊我的民以色列，如密雲遮蓋地面。\textcolor{red}{末後的日子}，我必帶你來攻擊我的地，到我在外邦人眼前，在你身上顯為聖的時候，好叫他們認識我。

\subsubsection{但以理書}
但以理書 2:28
只有一位在天上的神能顯明奧祕的事，他已將\textcolor{red}{日後}必有的事指示尼布甲尼撒王。你的夢和你在床上腦中的異象是這樣：

但以理書 10:14
現在我來，要使你明白本國之民\textcolor{red}{日後}必遭遇的事，因為這異象關乎後來許多的日子。」

\subsubsection{何西阿書}
何西阿書 3:5
後來以色列人必歸回，尋求他們的神耶和華和他們的王大衛。在\textcolor{red}{末後的日子}，必以敬畏的心歸向耶和華，領受他的恩惠

\subsubsection{约珥书}
使徒行傳 2:17
『神說：在\textcolor{red}{末後的日子}，我要將我的靈澆灌凡有血氣的，你們的兒女要說預言，你們的少年人要見異象，老年人要做異夢。


\subsubsection{彌迦書 4:1-3}
彌迦書 4:1-3
\textcolor{red}{末後的日子}，耶和華殿的山必堅立超乎諸山，高舉過於萬嶺，萬民都要流歸這山。

\subsubsection{撒迦利亞書}
撒迦利亞書 8:22
必有列邦的人和強國的民來到耶路撒冷，尋求萬軍之耶和華，懇求耶和華的恩
撒迦利亞書 10:8-12
「我要發嘶聲，聚集他們，因我已經救贖他們。他們的人數必加增，如從前加增一樣。
撒迦利亞書 14:9,20
耶和華必做全地的王，\textcolor{red}{那日}耶和華必為獨一無二的，他的名也是獨一無二的。

\subsubsection{提摩太後書 3:1}
你該知道，\textcolor{red}{末世}必有危險的日子來到。

\subsubsection{希伯來書 1:2}
就在這\textcolor{red}{末世}藉著他兒子曉諭我們；又早已立他為承受萬有的，也曾藉著他創造諸世界。

\subsubsection{彼得後書 3:3}
第一要緊的，該知道在\textcolor{red}{末世}必有好譏誚的人隨從自己的私慾出來譏誚說：

\section{以賽亞書}
\subsection{summery}
\subsubsection{summery 1}
\begin{itemize}
\item 一、不要信靠自己（1~12章） 
\item 二、不要信靠列国的帮助（13~23章） 
\item 三、信靠全地的主（24~27章） 
\item 四、信靠　神的要素（28~33章） 
\item 五、信靠　神的盼望（34~35章） 
\item 六、希西家作为信靠　神的榜样（36~39章） 
\item 七、不要信靠偶像（40~48章） 
\item 八、信靠义仆（49~55章） 
\item 九、信靠　神的人的胜利（56~66章） 
\end{itemize}

\subsubsection{summery 2}
\begin{itemize}
\item I. 亚述时期——争战及胜利 (1-35 章)
\begin{itemize}
\item A. 关于犹大和耶路撒冷的预言(1-12 章)
\item B. 关于列国的预言(13-27 章)
\item C. 真正的救赎之来源(28-35 章)
\end{itemize}
\item II. 历史间歇(36-39 章)
\item III. 巴比伦时期——苦难时期的盼望(40-66 章)
\begin{itemize}
\item A. 唯一真神与诸多偶像(40-48 章)
\item B. 通过受苦的仆人而得救(49-53 章)
\item C. 神子民将来的荣耀(54-66 章) 
\end{itemize}
\end{itemize}

\subsubsection{summery 3}
\begin{itemize}
\item I. 亚述时期(向著将要被擄之民而言,論亞述)——責備 (1-35 章)
\begin{itemize}
\item A. 責備猶大(1-12 章)
\begin{itemize}
\item 1. 猶大國遭審判(1-6 章)
\item 2. 責備警誡與預言(7-12 章)
\end{itemize}
\item B. 十一國遭宣判(13-23 章)
\item C. 責備猶大(24-34 章)
\begin{itemize}
\item 1. 預言猶大之將來(24-27 章)
\item 2. 警戒猶大之現在(28-35 章)
\end{itemize}
\end{itemize}
\item II. 過渡时期 (36-39 章)
\item III. 巴比伦时期(向被擄後之民而言,論巴比倫)(40-66 章)
\begin{itemize}
\item A. 回歸的應許（神是以色列人的救主）(40-48 章)
\item B. 拯救的應許（神是外邦人的救主）(49-57 章)
\item C. 復興的應許（神是以色列人的榮耀）(58-66 章) 
\end{itemize}
\end{itemize}


\subsection{彌賽亞（基督）生平的預言}
\begin{table}[H]
\centering
\begin{tabular}{p{1.5cm}|p{10.cm}}\hline\hline
\multicolumn{1}{c|}{彌賽亞（基督）生平}&\multicolumn{1}{c}{以賽亞書經文}\\\hline
出生&7:14，9:6-7\\\hline
身份&11:1，10\\\hline
名字&7:14，9:6，8:8，32:1，42:1\\\hline
傳道&11:2，42:1，61:1 – 3\\\hline
工作&11:3 – 4，35: 5 – 6，50:5 – 6，61:1 – 3\\\hline
品德&42:1 – 4\\\hline
受死&52:13 – 53:12，59: 20\\\hline
復活& 25: 8，53:10 – 11\\\hline
再臨& 63:1 – 6，64:1\\\hline
作王& 2: 2 – 4，11: 3 – 16，32:1，33:17 – 24，65: 12 –19\\\hline
\end{tabular}
\end{table}

\subsection{新約引用以賽亞書之處}
新約使徒多次引用以賽亞書為耶穌作見證，除詩篇外以賽亞書被引用為最多。
\begin{table}[H]
\centering
\begin{tabular}{p{3.5cm}|p{13.cm}}\hline\hline
\multicolumn{1}{c|}{引用者}&\multicolumn{1}{c}{經文}\\\hline
以賽亞書/詩篇&\\\hline
以賽亞書/彌迦書&2:1-3/4:1-3 \\\hline
以賽亞書/馬太&9:1-2/4:14-16, 8:17, 12:17-21, 9:1-2,53: 4,42: 1 – 4\\\hline
以賽亞書/馬可&15:28 53:12\\\hline
以賽亞書/路加&22:37 53:12/4:16 - 21 61:1 – 2(耶穌)\\\hline
以賽亞書/約翰&12:38 – 41 53:1, 6:9 – 10/1:23 40:3(施洗約翰)\\\hline
以賽亞書/使徒行傳&8:32-33 53:7-8/8:28-33 53:7-8(埃提阿伯太監)/28:25-27 6:9-10(保羅 )\\\hline
以賽亞書/彼得前書&2:22 – 25 53:5 – 6\\\hline
以賽亞書/羅馬書&1:9/9:27-28  9:29 10:11 10:15 10:16 10:20 – 21 15:12 10:22 – 23 1:9 49:23 52:7 53:1 65:1 – 2 11:10 \\\hline
\end{tabular}
\end{table}


%\chapter{启示录总结}

%\section{图表}
\subsection{以賽亞書预言}
%\begin{changemargin}{-1cm}{-1cm}
\begin{sidewaystable} % <-- HERE
\begin{tikzpicture}[thick,snake=zigzag, line before snake = 5mm, line after snake = 5mm]
   \draw (-2.5,5.2) node[below=3pt] {} node[above=3pt] {乌西雅};  
  \draw (-2.5,4.8) node[below=3pt] {} node[above=3pt] {死};  
    \draw (-2.5,4.4) node[below=3pt] {} node[above=3pt] {以};  
  \draw (-2.5,4.) node[below=3pt] {} node[above=3pt] {赛};  
     \draw (-2.5,3.6) node[below=3pt] {} node[above=3pt] {亚};   
    \draw (-2.5,3.2) node[below=3pt] {} node[above=3pt] {见};   
     \draw (-2.5,2.8) node[below=3pt] {} node[above=3pt] {主};   
    \draw (-2.5,2.4) node[below=3pt] {} node[above=3pt] {蒙}; 
    \draw (-2.5,2.) node[below=3pt] {} node[above=3pt] {召}; 
    \draw (-2.5,1.4) node[below=3pt] {} node[above=3pt] {740}; 
   \draw (-2.5,1.) node[below=3pt] {} node[above=3pt] {BC}; 
  

  \draw (-1.75,4.4) node[below=3pt] {} node[above=3pt] {亚};  
  \draw (-1.75,4.) node[below=3pt] {} node[above=3pt] {述};  
     \draw (-1.75,3.6) node[below=3pt] {} node[above=3pt] {攻};   
    \draw (-1.75,3.2) node[below=3pt] {} node[above=3pt] {打};   
     \draw (-1.75,2.8) node[below=3pt] {} node[above=3pt] {亚};   
    \draw (-1.75,2.4) node[below=3pt] {} node[above=3pt] {哈}; 
    \draw (-1.75,2.) node[below=3pt] {} node[above=3pt] {斯}; 
    \draw (-1.75,1.4) node[below=3pt] {} node[above=3pt] {734}; 
   \draw (-1.75,1.) node[below=3pt] {} node[above=3pt] {BC}; 
                  
   %  \draw (0.5,2) node[below=3pt] {} node[above=3pt] {};                         
   %  \draw (2.,2) node[below=3pt] {} node[above=3pt] {};                              
     \draw [dashed] (0.5,4.5) -- (0.5,7.8);
 \draw [>=triangle 45, ->] (0.5,7.8) -- (3.5,7.8);
   \draw (1.9,7.8) node[below=3pt] {} node[above=3pt] {以赛亚预言的事件};       
 
   %  \draw (0.5,4.) node[below=3pt] {} node[above=3pt] {亚};   
    \draw (-1,3.6) node[below=3pt] {} node[above=3pt] {亚};  
     \draw (-1,3.2) node[below=3pt] {} node[above=3pt] {述};   
    \draw (-1,2.8) node[below=3pt] {} node[above=3pt] {灭};   
     \draw (-1,2.4) node[below=3pt] {} node[above=3pt] {亚};   
    \draw (-1,2.) node[below=3pt] {} node[above=3pt] {兰}; 
     \draw (-1,1.4) node[below=3pt] {} node[above=3pt] {732}; 
   \draw (-1,1.) node[below=3pt] {} node[above=3pt] {BC}; 
  
    \draw (-.25,4.) node[below=3pt] {} node[above=3pt] {亚};   
    \draw (-.25,3.6) node[below=3pt] {} node[above=3pt] {述};  
     \draw (-.25,3.2) node[below=3pt] {} node[above=3pt] {掳};   
    \draw (-.25,2.8) node[below=3pt] {} node[above=3pt] {以};   
     \draw (-.25,2.4) node[below=3pt] {} node[above=3pt] {色};   
    \draw (-.25,2.) node[below=3pt] {} node[above=3pt] {列}; 
     \draw (-.25,1.4) node[below=3pt] {} node[above=3pt] {722}; 
   \draw (-.25,1.) node[below=3pt] {} node[above=3pt] {BC}; 

   \draw (.5,4.4) node[below=3pt] {} node[above=3pt] {亚};   
     \draw (.5,4.) node[below=3pt] {} node[above=3pt] {述};   
    \draw (.5,3.6) node[below=3pt] {} node[above=3pt] {攻};  
     \draw (.5,3.2) node[below=3pt] {} node[above=3pt] {打};   
    \draw (.5,2.8) node[below=3pt] {} node[above=3pt] {希};   
     \draw (.5,2.4) node[below=3pt] {} node[above=3pt] {西};   
    \draw (.5,2.) node[below=3pt] {} node[above=3pt] {家}; 
     \draw (.5,1.4) node[below=3pt] {} node[above=3pt] {722}; 
   \draw (.5,1.) node[below=3pt] {} node[above=3pt] {BC}; 


    \draw (1.25,4.4) node[below=3pt] {} node[above=3pt] {玛};   
     \draw (1.25,4.) node[below=3pt] {} node[above=3pt] {拿};   
    \draw (1.25,3.6) node[below=3pt] {} node[above=3pt] {西};  
     \draw (1.25,3.2) node[below=3pt] {} node[above=3pt] {杀};   
    \draw (1.25,2.8) node[below=3pt] {} node[above=3pt] {以赛亚};   
     \draw (1.25,2.4) node[below=3pt] {} node[above=3pt] {西};   
    \draw (1.25,2.) node[below=3pt] {} node[above=3pt] {家}; 
     \draw (1.25,1.4) node[below=3pt] {} node[above=3pt] {722}; 
   \draw (1.25,1.) node[below=3pt] {} node[above=3pt] {BC}; 

    \draw [dashed] (-1.2,5) -- (-1.2,5.6);
     \draw [dashed] (2,1.0) -- (2,5.6);
   \draw [>=triangle 45, <->] (-1.2,5.6) -- (2,5.6);
    \draw (2,.4) node[below=3pt] {} node[above=3pt] {680}; 
   \draw (2,6.9) node[below=3pt] {} node[above=3pt] {以}; 
   \draw (2,6.6) node[below=3pt] {} node[above=3pt] {赛}; 
   \draw (2,6.3) node[below=3pt] {} node[above=3pt] {亚}; 
   \draw (2,6) node[below=3pt] {} node[above=3pt] {死}; 
\draw[>=triangle 45, ->] (2,6.2) -- (2,5.56);
    \draw (2,0) node[below=3pt] {} node[above=3pt] {BC}; 
    \draw (.3,5.56) node[below=3pt] {} node[above=3pt] {以赛亚服事>60年};   
  
     \draw (2.5,4.8) node[below=3pt] {} node[above=3pt] {南};                                
     \draw (2.5,4.4) node[below=3pt] {} node[above=3pt] {国};   
    \draw (2.5,4.) node[below=3pt] {} node[above=3pt] {1次};                                
     \draw (2.5,3.6) node[below=3pt] {} node[above=3pt] {被};   
    \draw (2.5,3.2) node[below=3pt] {} node[above=3pt] {巴};                                
     \draw (2.5,2.8) node[below=3pt] {} node[above=3pt] {比};   
    \draw (2.5,2.4) node[below=3pt] {} node[above=3pt] {伦};                                    
    \draw (2.5,2.) node[below=3pt] {} node[above=3pt] {掳};                                                   
     \draw (2.5,1.4) node[below=3pt] {} node[above=3pt] {606}; 
   \draw (2.5,1.) node[below=3pt] {} node[above=3pt] {BC}; 


   \draw (3.5,4.8) node[below=3pt] {} node[above=3pt] {南};   
    \draw (3.5,4.4) node[below=3pt] {} node[above=3pt] {国};   
    \draw (3.5,4.) node[below=3pt] {} node[above=3pt] {3次};                                
     \draw (3.5,3.6) node[below=3pt] {} node[above=3pt] {被};   
    \draw (3.5,3.2) node[below=3pt] {} node[above=3pt] {巴};                                
     \draw (3.5,2.8) node[below=3pt] {} node[above=3pt] {比};   
    \draw (3.5,2.4) node[below=3pt] {} node[above=3pt] {伦};                                    
    \draw (3.5,2.) node[below=3pt] {} node[above=3pt] {掳};                                                   
     \draw (3.5,1.4) node[below=3pt] {} node[above=3pt] {586}; 
   \draw (3.5,1.) node[below=3pt] {} node[above=3pt] {BC}; 


   \draw (4.5,3.2) node[below=3pt] {} node[above=3pt] {南};                                         
   \draw (4.5,2.8) node[below=3pt] {} node[above=3pt] {国};                                            
   \draw (4.5,2.4) node[below=3pt] {} node[above=3pt] {归};                                       
   \draw (4.5,2.0) node[below=3pt] {} node[above=3pt] {回};                                            
    \draw (4.5,1.4) node[below=3pt] {} node[above=3pt] {536}; 
   \draw (4.5,1.) node[below=3pt] {} node[above=3pt] {BC}; 
   \draw[>=triangle 45, <->] (2.5,5.6) -- (4.5,5.6);
   \draw (2.5,5.4) -- (2.5,5.6);
  \draw (4.5,3.8) -- (4.5,5.6);
 \draw (3.5,5.6) node[below=3pt] {} node[above=3pt] {被掳70年};                                         

   \draw (6,3.6) node[below=3pt] {} node[above=3pt] {弥};                                         
   \draw (6,3.2) node[below=3pt] {} node[above=3pt] {赛};                                         
   \draw (6,2.8) node[below=3pt] {} node[above=3pt] {亚};                                            
   \draw (6,2.4) node[below=3pt] {} node[above=3pt] {一次};                                          
  \draw (6,2.0) node[below=3pt] {} node[above=3pt] {来};                                                
   \draw (6,1.35) node[below=3pt] {} node[above=3pt] {3}; 
   \draw (6,1.) node[below=3pt] {} node[above=3pt] {BC}; ;                                         

   \draw (7,3.2) node[below=3pt] {} node[above=3pt] {耶};                                         
   \draw (7,2.8) node[below=3pt] {} node[above=3pt] {稣};                                         
   \draw (7,2.4) node[below=3pt] {} node[above=3pt] {工};                                            
   \draw (7,2.) node[below=3pt] {} node[above=3pt] {作};                                          
 % \draw (7,1.6) node[below=3pt] {} node[above=3pt] {死};                                                

   \draw (8,3.2) node[below=3pt] {} node[above=3pt] {耶};                                         
   \draw (8,2.8) node[below=3pt] {} node[above=3pt] {稣};                                         
   \draw (8,2.4) node[below=3pt] {} node[above=3pt] {被};                                            
   \draw (8,2.) node[below=3pt] {} node[above=3pt] {钉};                                          
 
  \draw (9,3.2) node[below=3pt] {} node[above=3pt] {耶};                                         
   \draw (9,2.8) node[below=3pt] {} node[above=3pt] {稣};                                         
   \draw (9,2.4) node[below=3pt] {} node[above=3pt] {复};                                            
   \draw (9,2.) node[below=3pt] {} node[above=3pt] {活};                                          
  \draw (9,1.35) node[below=3pt] {} node[above=3pt] {30}; 
   \draw (9,1.) node[below=3pt] {} node[above=3pt] {AD}; ;                                         


   \draw (10,3.2) node[below=3pt] {} node[above=3pt] {圣};                                         
   \draw (10,2.8) node[below=3pt] {} node[above=3pt] {灵};                                         
   \draw (10,2.4) node[below=3pt] {} node[above=3pt] {降};                                            
   \draw (10,2.) node[below=3pt] {} node[above=3pt] {临};                                          
  \draw (10,1.35) node[below=3pt] {} node[above=3pt] {30}; 
   \draw (10,1.) node[below=3pt] {} node[above=3pt] {AD}; ;                                         
\draw[>=triangle 45, <->] (9,5) -- (10,5);
 \draw (9.5,5) node[below=3pt] {} node[above=3pt] {40天};                                         
 \draw (9,3.8) -- (9,5); 
\draw (10,3.8) -- (10,5);


   \draw (12,3.2) node[below=3pt] {} node[above=3pt] {教};                                         
   \draw (12,2.8) node[below=3pt] {} node[above=3pt] {会};                                            
   \draw (12,2.4) node[below=3pt] {} node[above=3pt] {被};                                          
   \draw (12,2.) node[below=3pt] {} node[above=3pt] {提};                                                
  \draw (12,1.35) node[below=3pt] {} node[above=3pt] {a}; 
\draw[>=triangle 45, ->] (12.5,0.8) -- (12.5,1.5);
\draw (12.5,0.2) node[below=3pt] {} node[text=blue,above=3pt] {3:18-26};
\draw (12.5,-0.1) node[below=3pt] {} node[text=blue,above=3pt] {4:1};

\draw[>=triangle 45, ->] (14.,0.8) -- (14.,1.5);
\draw (14.,0.2) node[below=3pt] {} node[text=blue,above=3pt] {2:13-21};
\draw (14.,-0.1) node[below=3pt] {} node[text=red,above=3pt] {26:20-21};

   \draw (15,3.6) node[below=3pt] {} node[above=3pt] {弥};                                         
   \draw (15,3.2) node[below=3pt] {} node[above=3pt] {赛};                                         
   \draw (15,2.8) node[below=3pt] {} node[above=3pt] {亚};                                            
   \draw (15,2.4) node[below=3pt] {} node[above=3pt] {二次};                                          
  \draw (15,2.) node[below=3pt] {} node[above=3pt] {来};                                                
 \draw (15,1.35) node[below=3pt] {} node[above=3pt] {b}; 

   \draw (16,3.6) node[below=3pt] {} node[above=3pt] {以色};                                         
   \draw (16,3.2) node[below=3pt] {} node[above=3pt] {列和};                                         
   \draw (16,2.8) node[below=3pt] {} node[above=3pt] {犹大};                                            
   \draw (16,2.4) node[below=3pt] {} node[above=3pt] {二次};                                          
  \draw (16,2.) node[below=3pt] {} node[above=3pt] {归回};                                                
 \draw (16,1.35) node[below=3pt] {} node[above=3pt] {c};

  \draw (17,3.2) node[below=3pt] {} node[above=3pt] {一};                                         
   \draw (17,2.8) node[below=3pt] {} node[above=3pt] {次};                                            
   \draw (17,2.4) node[below=3pt] {} node[above=3pt] {复};                                          
  \draw (17,2.) node[below=3pt] {} node[above=3pt] {活};                                                
 \draw (17,1.35) node[below=3pt] {} node[above=3pt] {d};
\draw[>=triangle 45, ->] (17,1) -- (17,1.5);
\draw (16.9,0.55) node[below=3pt] {} node[text=blue,above=3pt] {26:19};



  \draw (18,3.2) node[below=3pt] {} node[above=3pt] {审判};                                         
   \draw (18,2.8) node[below=3pt] {} node[above=3pt] {以色};                                            
   \draw (18,2.4) node[below=3pt] {} node[above=3pt] {列和};                                          
  \draw (18,2.) node[below=3pt] {} node[above=3pt] {国列};                                                
 \draw (18,1.35) node[below=3pt] {} node[above=3pt] {e};
\draw[>=triangle 45, ->] (18,0.8) -- (18,1.5);
\draw (18,0.2) node[below=3pt] {} node[text=blue,above=3pt] {1:24-31};


 \draw (19,3.6) node[below=3pt] {} node[above=3pt] {撒};                                                
  \draw (19,3.2) node[below=3pt] {} node[above=3pt] {旦};                                                
  \draw (19,2.8) node[below=3pt] {} node[above=3pt] {被};                                                   \draw (19,2.4) node[below=3pt] {} node[above=3pt] {困};                                                    \draw (19,2.) node[below=3pt] {} node[above=3pt] {绑};                                                
 \draw (19,1.35) node[below=3pt] {} node[above=3pt] {f};
\draw[>=triangle 45, ->] (19,0.8) -- (19,1.5);
\draw (19.25,0.2) node[below=3pt] {} node[text=blue,above=3pt] {27:1};


 \draw (20,3.6) node[below=3pt] {} node[above=3pt] {千};                                                
  \draw (20,3.2) node[below=3pt] {} node[above=3pt] {禧};                                                
  \draw (20,2.8) node[below=3pt] {} node[above=3pt] {年};                                                   \draw (20,2.4) node[below=3pt] {} node[above=3pt] {开};                                                    \draw (20,2.) node[below=3pt] {} node[above=3pt] {始};                                                
 \draw (20,1.35) node[below=3pt] {} node[above=3pt] {g};

\draw[>=triangle 45, ->] (20.75,0.8) -- (20.75,2);
\draw (20.75,0.2) node[below=3pt] {} node[text=blue,above=3pt] {2:1-5};
\draw (20.75,-0.1) node[below=3pt] {} node[text=blue,above=3pt] {4:2-6};

  \draw (21.5,4.8) node[below=3pt] {} node[above=3pt] {撒旦};                                                
  \draw (21.5,4.4) node[below=3pt] {} node[above=3pt] {出};                                                   \draw (21.5,4.) node[below=3pt] {} node[above=3pt] {无底坑};                                                    \draw (21.5,3.6) node[below=3pt] {} node[above=3pt] {诱惑};                                                
\draw (21.5,3.2) node[below=3pt] {} node[above=3pt] {列国};                                               
\draw (21.5,2.8) node[below=3pt] {} node[above=3pt] {围攻};                                                    \draw (21.5,2.4) node[below=3pt] {} node[above=3pt] {耶路};                                                
\draw (21.5,2.) node[below=3pt] {} node[above=3pt] {撒冷};                                               
 \draw (21.5,1.35) node[below=3pt] {} node[above=3pt] {h};

\draw (22.5,4.8) node[below=3pt] {} node[above=3pt] {所有};                                                
\draw (22.5,4.4) node[below=3pt] {} node[above=3pt] {的人};                                                   \draw (22.5,4.) node[below=3pt] {} node[above=3pt] {复活};                                                   
\draw (22.5,3.6) node[below=3pt] {} node[above=3pt] {白色};                                                
\draw (22.5,3.2) node[below=3pt] {} node[above=3pt] {大};                                                   \draw (22.5,2.8) node[below=3pt] {} node[above=3pt] {宝};                                                    \draw (22.5,2.4) node[below=3pt] {} node[above=3pt] {座};                                                
\draw (22.5,2.0) node[below=3pt] {} node[above=3pt] {审判};                                                
 \draw (22.5,1.35) node[below=3pt] {} node[above=3pt] {i};

  \draw (23.5,3.2) node[below=3pt] {} node[above=3pt] {进};                                                
  \draw (23.5,2.8) node[below=3pt] {} node[above=3pt] {入};                                                
    \draw (23.5,2.4) node[below=3pt] {} node[above=3pt] {永};                                                          
   \draw (23.5,2.) node[below=3pt] {} node[above=3pt] {世};                                                          
\draw (23.5,1.35) node[below=3pt] {} node[above=3pt] {j};
 %
      \draw (10,3.8) -- (10,5); 
      \draw (12,3.8) -- (12,5); 
\draw[>=triangle 45, <->] (10,5) -- (12,5); 
  \draw (10.5,5.0) node[below=3pt] {} node[above=3pt] {教};                                                
  \draw (10.9,5.0) node[below=3pt] {} node[above=3pt] {会};                                                
    \draw (11.3,5.0) node[below=3pt] {} node[above=3pt] {时};                                                          
   \draw (11.7,5.0) node[below=3pt] {} node[above=3pt] {期};                                                          


      \draw (15,4) -- (15,5); 
\draw[>=triangle 45, <->] (12,5) -- (15,5); 
  \draw (13.,5.0) node[below=3pt] {} node[above=3pt] {圣};                                                
  \draw (13.4,5.0) node[below=3pt] {} node[above=3pt] {徒};                                                
   \draw (13.8,5.0) node[below=3pt] {} node[above=3pt] {得};                                                          
   \draw (14.2,5.0) node[below=3pt] {} node[above=3pt] {赏};                                                          
   \draw (13.,4.36) node[below=3pt] {} node[above=3pt] {地};                                                
   \draw (13.4,4.36) node[below=3pt] {} node[above=3pt] {上};                                                
    \draw (13.8,4.36) node[below=3pt] {} node[above=3pt] {灾};                                                          
   \draw (14.2,4.36) node[below=3pt] {} node[above=3pt] {难};                                                          


       \draw (-2.5,2) -- (24.5,2);  
        \draw (0.5,2.1) -- (0.5,2);       
        \draw (1.5,2.1) -- (1.5,2);    
        \draw (3.5,2.1) -- (3.5,2);   
       \draw (16,2.1) -- (16.,2);     
       \draw (17.,2.1) -- (17.,2); 
     \draw (18.,2.1) -- (18.,2); 
       \draw (19.,2.1) -- (19.,2);           
       \draw (21.5,2.1) -- (21.5,2); 
      \draw (22.5,2.1) -- (22.5,2);  
      \draw (23.5,2.1) -- (23.5,2);  
                                                          
      \fill (2.5,2)  circle[radius=3pt];
   \fill (6.,2)  circle[radius=3pt]; 
   \fill (15.,2)  circle[radius=3pt];                                           
     \fill (20,2)  circle[radius=3pt]; 
     \fill (23.5,2)  circle[radius=3pt];  
\end{tikzpicture}
\end{sidewaystable} 
%\end{changemargin}





























